\begin{statementbox}
\textbf{条件与变量。} $\triangle ABC$ 为非退化三角形，记
$AB=p$、$AC=q$、$BC=a$。$A$ 处外角平分线所在直线记为 $\ell$。
若 $p\ne q$，则 $\ell$ 与直线 $BC$ 有唯一交点 $D$，且 $D$ 在线段 $BC$ 外部。
以下 $BD,DC,AB,AC$ 均指\textbf{正的普通长度}。

\textbf{外角平分线比例与外分关系：}
\[
 \boxed{\frac{BD}{DC}=\frac{AB}{AC}=\frac pq},\qquad
 \boxed{|BD-DC|=BC=a}.
\]
比例中的 $B$ 对应邻边 $AB$，$C$ 对应邻边 $AC$；不能因是外角就取倒数。

\textbf{先由邻边判断外分点在哪一侧：}
\[
\begin{array}{c|c|c}
 \text{邻边比较}&\text{直线上的顺序}&\text{长度关系}\\ \hline
 p>q&B,C,D&BD=BC+DC\\
 p<q&D,B,C&DC=BC+BD
\end{array}
\]
联立比例与差，可直接得到
\[
 BD=\frac{ap}{|p-q|},\qquad DC=\frac{aq}{|p-q|}.
\]
若取 $B$ 的坐标为 $0$、$C$ 的坐标为 $a$，则
\[
 x_D=\frac{ap}{p-q}.
\]
正负号在这里表示点的方向，不表示普通长度可以为负。

\textbf{等腰边界。} 若 $p=q$，则 $\ell\parallel BC$，\textbf{没有有限交点 $D$}，不能把 $p=q$ 代入上面的分式。

\textbf{与内角平分线对照。} 内角平分线交 $BC$ 于 $E$ 时，
$BE/EC=p/q$，但 $BE+EC=a$；内分用和，外分用差。
\end{statementbox}
